\documentclass[12pt]{article}
\usepackage[utf8]{inputenc}
\usepackage[T1]{fontenc}
\usepackage{hyperref}
\title{Sawknah Language Resources for K-12 Curriculum Development}
\author{Compiled by Claude Sonnet 5.5 (claude-sonnet-5-5) for Andr\'as Kornai}
\date{}
\begin{document}
\maketitle

\section{Introduction}
Sawknah, also called Sokna (ISO 639-3 swn, Glottocode sawk1238), is a presumably extinct Eastern Berber (Afro-Asiatic) language once spoken in the town of Sokna (Isuknan) and the village of Fuqaha in northeastern Fezzan, Libya. Wikipedia reports that its main documentation recorded only a handful of old speakers. Endangered Languages lists it as dormant and notes a non-fluent speaker's wordlist posted online in 2010. Lameen Souag's 2016 article on vocabularies from 1850 reports only five fluent speakers in 1915. There is no school tradition. Any programme would be one of heritage awareness and revitalization from documentation.

\section{Primary Educational Resources}
None found. No primers, textbooks or school programmes exist in the search results.

\section{Linguistic and Reference Materials}
\begin{itemize}
\item Wikipedia article on the Sokna language: \url{https://en.wikipedia.org/wiki/Sokna_language}
\item Wikipedia article on Awjila, the closely related living Eastern Berber language of Cyrenaica, useful for comparison: \url{https://en.wikipedia.org/wiki/Awjila_language}
\item Wiktionary category for Sokna language entries: \url{https://en.wiktionary.org/wiki/Category:Sokna_language}
\item Endangered Languages Project entry: \url{https://endangeredlanguages.com/lang/294?hl=en}
\item Numerals page from the Max Planck EVA channumerals database: \url{https://lingweb.eva.mpg.de/channumerals/Sawknah.htm}
\end{itemize}

\section{Religious and Community Resources}
None found.

\section{Digital and Audio Resources}
No audio or video resources were found in the search results. The sources are wordlists and published linguistic descriptions; the original documentation should be consulted through the bibliography cited on Wikipedia.

\section{Curriculum Development Context}
A heritage curriculum could introduce the history and place of Sokna and Fezzan, the Eastern Berber family and the available wordlists, with comparison to Awjila. Learners should be told clearly that the language is poorly documented and that reconstruction would rest on limited sources.

\end{document}